% Source scan 673; block decisions in selection.json

\quellseite{673}

a Frühe) das ist, er sagete ihnen zeitlich zuvor, daß sie Zeit genug hätten, allemal sich zu bessern, ehe die Straffe käme.

\section{Summarium}

Ist eine Historia, wie die Kinder des frommen Josia ungerathen und abgöttisch gewesen sind, und sie Gott darum gestraft habe, bis endlich der König zu Babel das Land gar verheeret und das Volck gefangen hinweg geführet hat, V. 17. 20. Solche Straffe soll uns ja ein Exempel und Warnung sein, daß wir uns für Abgötterei hüten, Gottes Wort mit Ernst aufnehmen und demselbigen nachwandeln sollen; wo wir anders nicht gleichen Zorn und Straffe mit den Juden tragen wollen, welche siebenzig ganzer Jahre unter den Heiden in allerlei Jammer und Elend, um solcher Sünden willen mußten leiden.

\part{Das Buch Esra}

\chapter{Das 1. Kapitel}

\section{Vorrede}

Bei diesem ersten Kapitel sollen wir merken zwei Stücke. Das erste ist eine Historia, wie das jüdische Volck wieder in sein Vaterland kommt aus dem babylonischen Gefängniß. Diese lehret, daß Gottes Verheißungen gewiß sein, er selbst wahrhaftig, und daß man göttlichem Wort billig und sicher glauben möge. Auch lehret die Geschichte, daß Gott allmächtig sey und deswegen die Seinen aus Creuz und Leiden erretten könne, auch alsdann, wenn es für der Vernunft und auch menschlichen Kräften unmöglich scheinet.

Das andere Stück ist das Exempel Kores und des babylonischen Volckes, daß, wie Kores des jüdischen Volckes: also die Obrigkeit der Kirchen Gottes sich erbarmen, ihrer sich annehmen, ihr Glück und Wohlfahrt fördern und fortsetzen soll; und, wie das jüdische Volck mit Silber und Gold geholfen hat zum Bau des Tempels zu Jerusalem: also soll jedermann Hülffe thun zur Förderung göttlichen Wortes und des rechten Gottesdienstes, und derwegen von Gott reiche Belohnung erwarten, und das soll geschehen, nicht mit Unwillen oder aus Zwang, sondern freiwillig. Denn wie die Schrift sagt: Einen fröhlichen Geber hat Gott lieb, 2 Cor. 9, 7. Und Sirach im 35. Capitel V. 20. spricht: Wer Gott dienet mit Lust, der ist ihm angenehm.

\section{Kores, oder Cyrus, entläßt die gefangenen Juden}
% Source scan 674; block decisions in selection.json

\quellseite{674}

\section{Summarium}

Bisher hat eure Liebe gehört, wie die Juden sich mit Abgötterei, falscher Lehre und Verachtung des Wortes Gottes in ihren Landen versündiget haben, und darum von den Babyloniern gefangen weggeführet sind, und siebenzig Jahr im Gefängniß blieben.

Diese Zeit aber ist der Prophet Daniel in einem großen Ansehen gewesen unter den Heiden. Und sind durch seine Predigt NebucadNezar und EvilMerodach bekehret worden. Aber unter dem dritten Könige Belsazer hat die Abgötterei und andere Untugend wieder gewachsen und überhand genommen.

Darum schicket Gott über Babylonien Darium aus Meden und Cyrum aus Persien; dieselben mußten eben an den Babyloniern die Verachtung des Wortes Gottes und andere Sünde strafen, wie sie vorhin die beiden Könige von Assyrien und Babel an den Juden gestrafet hatten. Solche beide Könige, Darius und Cyrus, sind durch den Propheten Daniel auch bekehret zum rechten Gott.

Weil aber Esaias und Jeremias von Cyrus hatten geweissagt, daß er sollte das Volk in ihr Land Juda ziehen und den Tempel wieder aufbauen lassen, nimmt Cyrus diesen Befehl an und folget solcher göttlichen Weissagung. Befiehlt auch dazu seinen Landen und Leuten (welches sonderlich ein fein fürstlich Werk ist), daß sie den Juden zu solchem Werk reiche Beisteuer thun und helfen sollen, V. 6. Also hat Gott seines Volks Gefängniß unter den Heiden dahin gebracht, daß sein Wort unter ihnen gepredigt und sie auch Gottes Volk und selig worden sind. Wie nun die Juden wieder zu Haus kommen, die Stadt Jerusalem und den Tempel wieder zu bauen angefangen und wie es ihnen darüber ergangen hat, werden folgende zwei Bücher, Esra und Nehemia anzeigen.

\section{Votum}

Der treue GOtt helfe, daß Obrigkeit und Unterthanen seiner Kirchen sich treulich annehmen, belohne ihnen alle ihre Wohlthaten reichlich, und steure allen Kirchfeinden durch seinen gewaltigen Arm, um der Ehre seines heiligen Namens willen, Amen.

\chapter{Das 2. Kapitel}

\section[Zahl und Register der wiederkehrenden Juden, und ihre]{Zahl und Register der wiederkehrenden Juden, und ihre Steuer}
% Source scan 675; block decisions in selection.json

% Scan 675: no apparatus retained; only biblical narrative/running matter.
% Source scan 676; block decisions in selection.json

\quellseite{676}

\chapter{Das 3. Kapitel}

\section{Vorrede}

Wir sind zwar von den mosaischen Opfern erlöst und befreiet, dieweil der HErr Christus sich nun selbst am Stamm des Kreuzes geopfert hat; dennoch, welche durch den Glauben des Leidens und Sterbens Christi theilhaftig worden sind, die sind gleichwohl schuldig, Opfer Gott zur Danksagung zu beweisen. Derwegen sollen wir Gott opfern einen betrübten, geängsteten Geist, von wegen unserer Sünden. Wir sollen Gott unsere Vernunft opfern und wider dieselbige glauben, was zu glauben ist, die gefangen nehmen und dem Gehorsam Christi unterwerfen. Wir sollen Gott Opfer bringen mit unserm Mund und Lippen, mit Beten und Danksagung. Wir sollen Gott opfern Gerechtigkeit gegen unsern Nächsten. Wir sollen Gott opfern Barmherzigkeit und die Werke der Liebe gegen den Menschen beweisen. Wir sollen Gott aufopfern unsere Kinder, durch Gebet, durch die Taufe, durch gute Unterweisung und Zucht und daß wir sie ihm willig und gern sollen folgen lassen, wenn er sie durch den Tod zu sich fordert. Wir sollen auch Gott unsern Leib aufopfern, mit Thun und Lassen, Leiden und Sterben und sollen kreuzigen unser Fleisch, sammt den bösen Lüsten und Begierden, und den alten Adam tödten. Dies betrachten wir auch bei dem 3. Kapitel.

\section[Vom Brandopferaltar. Fest der Laubhütten und Gründung des]{Vom Brandopferaltar. Fest der Laubhütten und Gründung des Tempels}
% Source scan 677; block decisions in selection.json

\quellseite{677}

\section{Summarium}

Ist eine Historia, wie Jesua und Serubabel endlich den Altar baueten, Es. 2., auf daß, ob man mit dem andern Gebäude so bald nicht möchte fertig werden, sie dennoch opfern und den Gottesdienst üben könnten. Denn weil Gott die Opfer selbst geboten und verordnet hatte, war es ein Werk des Gehorsams. Siehet man, wie ein fein biblisch Werk das sei, dazu helfen und rathen, daß der rechte Gottesdienst, so gefallen ist, wieder angerichtet und gefördert werde.

\section{Votum}

GOtt sei Lob und Dank, der zu aller Zeit christliche Obrigkeit gibt und erhält; der segne und bewahre ferner unsere liebe Obrigkeit und alle andere christliche Regenten, die ihn ehren und seine Kirche fördern, und auf ihrer Unterthanen Heil und Wohlfahrt trachten, Amen.

\chapter{Das 4. Kapitel}

\section{Vorrede}

Daß die gläubigen Bekenner und Diener des HErrn Christi angefeindet, fälschlich angegeben, verhindert und verfolget werden, ist nicht heuer neu. Denn solches hat sich bald nach der Erschaffung und Fall des Menschen angefangen, und währet noch also und wird bleiben bis ans Ende der Welt. Aber solches sollen wir Christen nicht achten, sondern sollen unter solchem Kreuz uns trösten, daß uns Gott dennoch lieben, erhalten und fördern will, und daß er uns reichlich wird erstatten, was uns die Welt Schaden und Nachtheil zufüget. Denn Gott läßt den Seinen hinweg nehmen die zeitlichen Güter, auf daß er ihnen gebe die ewigen himmlischen Güter und Reichthümer. Gott läßt den Seinen nehmen und abziehen ein geringes Kleid, damit er sie ziere und schmücke mit dem güldenen Kleid der ewigen Herrlichkeit. Gott läßt den Seinen das zeitliche Leben nehmen, daß er ihnen gebe das ewige Leben. Das läßt auch unsern Trost sein unter allem Haß, Neid, Widerwärtigkeit und Verfolgung der Welt und ihrer Helfer. Und höret darauf an das 4. Kapitel.

\section{Hinderung des Tempelbaues, von der Juda Feinden angestellet}
% Source scan 678; block decisions in selection.json

\quellseite{678}

\section{Summarium}

Dies ist eine Historia, in welcher man siehet, daß der Teufel dem rechten Gottesdienst feind ist, denselbigen nicht leiden kann, sondern mit Lügen und wie er kann, begehrt zu dämpfen. Denn das Volk in dem Lande durch den Teufel getrieben, hindert nicht allein den Bau des Tempels, V. 4, sondern verklagt auch fälschlich die Juden als Aufrührerische, daß sie in Verdacht bei dem Könige kämen, wenn sie die Stadt befestiget hätten, daß sie alsdann abfallen und andere mit ihnen zum Abfall bereden und sich stärken würden.
% Source scan 679; block decisions in selection.json

\quellseite{679}
Solcher Nachrede und Lügen müssen die Christen noch heutiges Tages gewarten und tragen um des Evangelii und des rechten Gottesdienstes willen. Denn alsbald sie davon nicht abstehen, noch etwas begeben wollen, wie sie denn ihres Gewissens und der Ehre Gottes halben nicht können, so heißt man sie ungehorsam, aufrührerische Leute. Aber es geschieht ihnen unrecht, wie den frommen Juden hie.

\section{Votum}

Der allmächtige Gott steure und wehre allen Feinden seiner Kirchen, stärke und erhalte uns in Verfolgung und unter allem Widerstand; und helfe, daß sein Wort bei uns geprediget und die rechten Gottesdienste geübet werden, bis an den jüngsten Tag. Amen.

\chapter{Das 5. Kapitel}

\section{Vorrede}

Der königliche Prophet David sagt im 55. Psalm V. 23.: Wirf dein Anliegen auf den Herrn, der wird dich versorgen und wird den Gerechten nicht ewiglich in Unruhe lassen. Hiermit lehret der heilige David: Kreuz, Trübsal, Angst und Noth soll und muß bei den Frommen sein: darein sollen sich auch alle Gottesfürchtige freiwillig und geduldig ergeben; sintemal es Gottes Wille ist, und weil wir des Herrn Christi Malzeichen und Sterben an unserm Leibe tragen, seines Leidens viel haben, mit Thränen säen, Angst und Leid haben müssen. Aber, wie Gott demüthiget, also erhöhet er auch, daß er tröstet durch sein kräftig Wort, und Kraft und Vermögen verleihet, daß wir die Noth überwinden können; und endlich hilft er auch unter und aus allem Kreuz, und thut solches, wenn es ihm wohlgefällt und nun die größte und äußerste Noth vorhanden ist. Dieser Lehre und Trost ist ein Exempel das 5. Kapitel.

\section{Fortsetzung des Tempelbaues}
% Source scan 680; block decisions in selection.json

\quellseite{680}

\section{Summarium}

Sie ist ein fein tröstlich Exempel, nachdem die Juden von den umliegenden Heiden sind verklagt und ihnen der Bau niedergelegt wird, will Gott sein Volk in solcher Anfechtung nicht gar ohne allen Trost lassen. Schicket zween Propheten, welche sie trösten und vermahnen sollten, an dem angefangenen Werk fortzufahren, V. 8. Item, er wendet des Landpflegers und anderer Amtleute Herz, daß sie dem Könige Dario anzeigen, daß ihnen von Cyro selbst zu bauen erlaubt sei. Also siehet man, daß Anfechtung und Widerwärtigkeit in göttlichen Sachen nicht außen bleiben und doch Gott allwege Hülfe und Trost schicket, und Mittel schaffet, daß sein Wille und nicht des Teufels, noch der argen Welt Wille geschehe.

\section{Votum}

Wir danken unserm treuen lieben Gott, daß er seine Kirche und Gemeine bisher so wunderlich erhält und fördert; der wolle an uns auch ferner diese seine Gnade beweisen, und allen verfolgten Christen Trost und Hülfe mittheilen und erzeigen, Amen.

\chapter{Das 6. Kapitel}

\section{Vorrede}

Wir sollen täglich und fleißig Gott anrufen, und beten für unsere Obrigkeit. Denn die Obrigkeit solches um uns verdienet und auch bedarf. Denn wie es bei Gott stehet und wir ihn darum bitten müssen, daß er guten fruchtbaren Regen der Erde geb, soll sie anders grünen und Früchte bringen: also muß er auch seinen Segen der Obrigkeit geben, soll sie nützlich und wohl regieren. Darum sollen wir ihn auch herzlich bitten: So hat es Gott auch befohlen, als dies St. Paulus 1 Tim. 2, 1. lehret, da er sagt: So ermahne ich, daß man für allen Dingen zuerst thue Bitte, Gebet, Fürbitte und Danksagung für alle Menschen, für die Könige und für alle Obrigkeit, auf daß wir ein geruhiges und stilles Leben führen mögen, in aller Gottseligkeit und Ehrbarkeit. Denn solches ist gut, dazu auch angenehm für Gott unsern Heiland. Ueber das sollen wir für die Obrigkeit beten, daß uns unser Gewissen nicht verklage. Welches geschehen würde, da wir der Obrigkeit die Ehre nicht thäten, die wir unsern Eltern schuldig sind, darunter die Fürbitte auch wird begriffen. Dieses lehret auch das 6. Kapitel.

\section{Vollführung und Einweihung des Tempels}
% Source scan 681; block decisions in selection.json

\quellseite{681}

\section{Summarium}

Ist eine schöne historia, wie der König Darius auf ein neues befehlet, man soll nicht allein die Juden an ihrem Bau ungehindert lassen, sondern auch von dem königlichen Einkommen dazu geben, was vonnöthen ist, V. 6. Solches ist ja eine feine fürstliche Tugend, wo die Obrigkeit sich des rechten Gottesdienstes annimmt und sich seiner Unkosten dauern läßt.

Sonderlich aber ist hie zu merken, daß der gottseelige König Darius begehrt, daß man soll zu Jerusalem für sein und seiner Kinder Leben bitten, und stellet eine harte Strafe und greulichen Fluch darauf, so jemand wider solchen seinen Befehl handeln, und die Juden hindern wollte.

\section{Votum}

Wir danken GOtt, der uns auch christliche Obrigkeit gibt und erhält; der segne, regiere, stärke und erhalte sie ferner, auf daß sie ihrem Regiment möge wohl fürstehen, daß er geehrt und die armen Unterthanen geschützt werden. Amen.

\chapter{Das 7. Capitel}

\section{Vorrede}

Die Obrigkeit soll ein Schutzherr sein beider Tafeln der heiligen Zehen Gebote, und neben der weltlichen Regierung die christliche Kirche treulich befördern. Darum wird dem römischen Kaiser in der Krönung nicht allein ein Scepter, sondern auch ein Reichsapfel, darauf ein Kreuz stehet, in die Hand gegeben, zur Anzeigung, daß die Regenten die heilsame Lehre von dem gekreuzigten Christo befördern, schützen und handhaben sollen: also hat ihm der Kaiser Carolus Magnus, wenn er in seiner Herrlichkeit herein gegangen, die Bibel und ein Schwert fürtragen lassen, zur Erinnerung, daß er und alle christliche Obrigkeit beide die christliche Kirche mit dem reinen Wort Gottes versorgen, und das Schwert im weltlichen Regiment zum Schutz der Frommen und zur Strafe der Bösen führen solle. So helfe der treue Gott, daß alle Obrigkeit ihm, als ihrem Herrn, diene mit Furcht, und daß sie küssen seinen Sohn, unsern Herrn Jesum Christum, daß er nicht zürne und sie umkommen auf dem Wege. Davon höret auch an das 7. Kapitel.

\section[Ezra bekommt Erlaubniß und Beförderung, den Gottesdienst zu]{Ezra bekommt Erlaubniß und Beförderung, den Gottesdienst zu bestellen}
% Source scan 682; block decisions in selection.json

% Scan 682: no apparatus retained; only biblical narrative/running matter.
% Source scan 683; block decisions in selection.json

\quellseite{683}

\section{Summarium}

Ist ein anderer schriftlicher Befehl, den der König Artahsastha Esra, dem Schriftgelehrten, gibt, V. 11., in welchem man stehet, wie dieser heidnische König so ein feines Herz und guten Willen zu dem Gott Israel und seinem Dienst hat, daß er nicht allein aus seiner Kammer verordnet, Silber und Gold darum zu kaufen, was zum Gottesdienst gehöret, sondern auch das Gewissen hat, wo er den Gottesdienst nicht also förderte, es möchte Gottes Zorn über sein Königreich und seine Kinder kommen, V. 23.

Es ist aber sonderlich hie zu merken, daß er die Priester und Leviten befreyet von allen Zinsen und Renten, daß man sie damit unbeschweret lasse, V. 24. Auch befiehlet er dem Esra, daß er sonderlich darauf soll Achtung geben, daß man rechte Prediger möge halten, und die, so noch nicht genugsam im Gesetz erfahren, recht unterweiset und gelehret würden, V. 25. Dies sind rechte fürstliche, ja bischöfliche Werke, dafür man billig (wie Esra hie thut, V. 27.) Gott danken soll, wo eine Obrigkeit sich dermaßen in die Sache schicket.

\section{Votum}

Wir danken GOtt herzlich auch für unsere liebe Obrigkeit; der regiere und erhalte sie, daß sie ferner die göttliche Lehre und treue Lehrer schütze und fördere, und belohne reichlich ihre väterliche Treue an ihnen selbst, an ihren Erben und Kindern und an ihren Hab und Gütern, Amen.

\chapter{Das 8. Capitel}

\section{Vorrede}

Die Vollträgigen und Trunkenen beten entweder gar nicht, oder beten ohne alle Andacht, und mit bösem Gewissen. Derwegen sollen wir uns für Fressen und Saufen hüten und ein nüchtern und mäßiges Leben führen, damit wir zum Gebet desto geschickter sein. Man schreibt vom römischen Regenten Romulus, daß derselbe in einer Collation wenig Weins getrunken, und da er gefragt ward, warum er so wenig getrunken? hat er darauf geantwortet: Er hätte eine wichtige Sache zu verrichten; also vielmehr sollen wir Christen, um des Gebets willen, daran uns trefflich viel gelegen, uns mäßig im Essen und Trinken halten. Daher setzet auch die Schrift gemeiniglich zusammen: Fasten und Beten. Als Matth. 7. 21. saget Christus: Dieses Geschlecht fähret nicht aus, denn durch Beten und Fasten. Und St. Petrus spricht 1. Epistel, 5. Kapitel V. 8.: Seid nüchtern und wachet. Und der alte Tobias sagt im 12. Kapitel V. 9.: Es ist ein köstlich Ding, das Gebet mit Fasten. Dieses lehret auch das 8. Kapitel.

\section{Erste Reise gen Jerusalem wird vollbracht}
% Source scan 684; block decisions in selection.json

\quellseite{684}

\section{Summarium}

Ist eine Historia, in welcher erzählet wird, welche mit Esra aus Babel gen Jerusalem gezogen seien. Und sonderlich ist hie zu merken, was es für eine Meinung habe, daß Esra, ehe er sich auf den Weg begibt, ein Fasten ausrufet, V. 21. Wie der König Josaphat 1 Chron. 20. auch thut, da die Moabiter ihm ins Land fielen. Denn wo eine gemeine Gefahr und Noth fürfiele, da geboten die Könige und Priester ein Fasten. Nicht wie die Papisten, daß sie damit Gott wollten dienen, sondern daß jedermann zum Gebet desto geschickter wäre, und sie sich also für Gott demüthigten und mit größerem Ernst bei Gott um Hülfe ansuchten. Denn wo ein voller Bauch ist, da wird wenig ernstliches Gebets sein. Darum Christus auch mehr denn einmal Fasten und Beten zusammen setzet. Von solchem Ernst des Gebets wissen unsere Papisten nichts, und wollen dennoch aus ihrem spöttlichen Fasten einen sonderlichen Gottesdienst machen.

\section{Votum}

Helfe der treue GOtt, daß wir der heiligen Exempel fleißig folgen, und GOtt Tag
% Source scan 685; block decisions in selection.json

\quellseite{685}
und Nacht und uns selbst dienen, mit nüchternem Leben und Gebet, und allen andern Tugenden, in Kraft des heiligen Geistes, Amen.

\chapter{Das 9. Kapitel}

\section{Vorrede}

Hieran ist alles gelegen, da wir selig werden wollen, daß uns Gott die Sünde vergebe, daß wir auch solches bitten und im Glauben annehmen müssen. Denn, so wir sagen: Wir haben keine Sünde! so verführen wir uns selbst, und die Wahrheit ist nicht in uns, sagt der heilige Apostel Johannes 1. Epistel im 1. Kapitel V. 8. So wird alle Sünde den Gläubigen vergeben um des Herrn Christ willen. Daher sagt St. Petrus in den Geschichten der Apostel im 10. Kapitel V. 43.: Von Christo zeugen alle Propheten, daß durch seinen Namen alle, die an ihn glauben, Vergebung der Sünde empfahen sollen. Denn Christus ist unser Erlöser, welcher bezahlet beide, den alten Rest der Erbsünde und alle andere wirkliche Sünde. Er hat uns eine gnädige Quitanz erlanget. Denn sein rosinfarben Blut ist durch unser Schuld-Buch durchgeflossen, zur Vergebung aller unserer Sünden. Es wird auch nimmermehr keine Rechnung in die himmlische Renterei eingenommen, wo sie nicht mit der schönen Rubrik des Blutes und Todes Jesu Christi zuvor cassirt und quittirt oder ausgelöscht ist. Davon lehret auch insgemein das 9. Kapitel.

\section{Etwas Beicht- und Bußgebet}

a Nagel und Zaun) ist geredt auf Sprichwortsweise, daß alles andere, Land und Leute umkommen sind, und sie, die Israeliten, noch übrig sind blieben, als ein Nagel vom Hause und ein Zaun vom Lande.
% Source scan 686; block decisions in selection.json

\quellseite{686}

\section{Summarium}

Dieß ist eine ernstliche Beichte und feines Gebet des frommen Esra, V. 6., daß GOtt die Sünde vergeben und nicht nach dem Verdienst strafen wolle, daß die Priester und Leviten heidnische Weiber genommen hatten, V. 12. Uns zum Exempel, daß wir nach gethaner Sünde dieselbigen für Gott von Herzen bekennen sollen, und uns nichts anders trösten, denn seiner Gnade und Barmherzigkeit und auf dieselbigen unser Gebet gründen.

\section{Votum}

Der treue barmherzige GOtt sei uns allen gnädig, vergebe uns unsere Sünde und lindere uns die Strafen, so wir mit unsern Sünden wohl verdienet haben; und helfe, daß wir uns allezeit seiner Gnade und unseres Herrn Christi Verdienstes mit Freuden trösten, Amen.

\chapter{Das 10. Kapitel}

\section{Vorrede}

Es gehören zu rechter Buße und Bekehrung zu Gott drei Stücke. Das erste ist herzliche Erkenntniß der Sünden und wahre Reue und Leid von wegen der Sünden. Davon lehret Gott im 66. Kapitel V. 2. Esaia also: Ich sehe an den Elenden und der zerbrochenes Geistes ist, und der sich fürchtet für meinem Wort. Das andere Stück der christlichen Buße ist der Glaube, daß Gottes Güte und Gnade viele tausend mal größer sei, denn aller Menschen Sünde und Uebertretung, und daß er allen, die ihn darum bitten, alle Sünden vergibt, um des Herrn Christi Verdienstes und Vorbitte willen. Dieses lehret David, Ps. 147, 11.: Der Herr hat Gefallen an denen, die ihn fürchten und auf seine Güte warten. Und St. Paulus sagt Col. 1, 14.: An Christo haben wir die Erlösung durch sein Blut, nämlich die Vergebung der Sünden. Das dritte Stück der wahren Buße ist neue Gehorsam nach Gottes Wort und Gebot, daß der Mensch von Sünde ablasse und in allen guten Werken sich übe, wie es ihm in seinem Stande gebühret, damit der siebenfältige Teufel nicht wieder einschleiche und das Letzte ärger werde, denn das Erste. Dies bestätiget St. Paulus an die Römer im 6. Kapitel V. 12.: Lasset die Sünde nicht herrschen in eurem sterblichen Leibe; ihr Gehorsam zu leisten in ihren Lüsten. Auch begebet nicht der Sünden eure Glieder zu Waffen der Ungerechtigkeit, sondern zu Waffen der Gerechtigkeit. Davon lehret auch das 10. und letzte Kapitel.

\section{Die fremden Weiber werden von den Israeliten abgeschaffet}
% Source scan 687; block decisions in selection.json

\quellseite{687}

\section{Summarium}

Es ist eine Historia, daß Esra daran ist, daß alle so heidnische Weiber genommen, ob sie gleich Kinder mit ihnen gezeuget hatten, dieselbigen sammt den Kindern von sich thun sollen, V. 3. Denn weil es wider Gottes Wort und Befehl geschehen war, konnte es keine rechte Ehe sein. Hie stehet man abermal, was rechtschaffene Buße ist, nämlich, daß man ihm die Sünde von Herzen lasse leid sein und davon ablasse, und hoffe darnach, Gott werde sie gnädiglich vergeben. Ohne Zweifel wird es wohl haben gethan, sich von Weib und Kind scheiden. Aber weil Got-
% Source scan 688; block decisions in selection.json

\quellseite{688}
tes Befehl da war, und ohne daß die Buße nicht konnte rechtschaffen sein, mußten sie es thun, oder zugleich Gottes Zorn und ein böses Gewissen tragen.

\section{Votum}

Helfe der gütige barmherzige GOtt, daß wir auch wahre Buße thun, in wahrer Buße alle Tage leben unser Lebenlang, an den Herrn Christum fest glauben, und unsern Glauben beständig beweisen mit guten Werken, gegen GOtt und gegen unsern Nächsten, nach seinen Geboten und Hülfe des heiligen Geistes, Amen.

\part{Das Buch Nehemia}

\chapter{Das 1. Kapitel}

\section{Vorrede}

Ein rechter Beter soll ein Mensch sein, der nicht nur habe Mund und Zunge, sondern auch heilige Hände, übersichtige Augen, ein gut Gewissen und ein gläubig Herz, das ist: ein rechter Beter soll sein Gebet nicht allein ein Mundwerk von Worten sein lassen, sondern er soll heilige Hände aufheben, ohne Zorn, Blut und Frevel. Des Beters Augen sollen nicht unter sich sehen, wie der Zweifler, die an Gott und seiner Gnade verzagen, sondern sollen über sich sehen gen Himmel, in fröhlicher Zuversicht auf den Namen Jesu Christi und göttliche Verheißung; der Beter soll haben ein gut Gewissen und gute Sache, daß die Person mit Gott versöhnet sei. Denn ein böses Gewissen kann nicht beten. Der Beter soll auch haben ein rechtschaffen Herz, brünstig in der Anrufung und gläubig auf Gottes Zusage und Verheißung. Denn wer also betet, daß er zweifelt an der Erhörung Gottes und allein sein Gebet dahin setzt auf Ebentheuer, der begehet zwei große Sünden: erstlich, daß er sein Gebet selbst zu Sünden und Schaden macht; zum andern, daß er seinen getreuesten und wahrhaftigsten Gott für einen Lügner acht, der seinem Wort und Zusage nicht könne oder wolle Folge thun. Derwegen wer wohl und nutzbarlich beten will, der schaue wohl zu, daß er ein rechter Beter sei. Dies betrachten wir auch bei dem 1. Kapitel.

\section{Nehemias Gebet für die Kinder Israel}
% Source scan 689; block decisions in selection.json

\quellseite{689}

\section{Summarium}

Dies ist ein schön Gebet, welches Nehemias thut, weil er der Juden Sache bei dem Könige handeln wollte, daß ihm Gott zu solchem Anbringen wollte Glück geben, V. 5. Und zum Exempel, daß wir nicht allein lernen, alle Dinge mit dem Gebet anfangen, auf daß wir desto mehr Glück haben, sondern daß wir auch sehen, wie ein recht Gebet müsse gestellet sein, nämlich, daß man sich erstlich für Gott demüthige und bekenne, daß wir Sünder und aller Strafe werth sind. Zum andern, daß man auf Gottes Verheißung und Zusage das Gebet stelle, wie Nehemias die Worte erzählet, in welchen Gott Mose zugesagt hat, er wolle auch der Sünder Gebet erhören und annehmen, wenn sie sich bekehren und sein Gebot halten.

\section{Votum}

Gott, der Brunnquell alles Guten, verleihe uns auch bußfertige, demüthige Herzen, und helfe, daß wir täglich fleißig beten, für unsere und der Christenheit gemeine Nothdurft, und daß wir in wahrem Glauben seiner väterlichen Hülfe erwarten, Amen.

\chapter{Das 2. Kapitel}

\section{Vorrede}

Der königliche Prophet David spricht im 145. Psalm V. 19.: Der HErr thut, was die Gottesfürchtigen begehren. Und der weise König Salomon in seinen Sprüchen sagt im 14. Kapitel V. 27.: Die Furcht des HErrn ist eine Quelle des Lebens, daß man meide die Stricke des Todes. Und Sirach im 1. Kapitel V. 18. spricht: Wer den HErrn fürchtet, dem wird es wohl gehen, und wenn er Trost bedarf, wird er gesegnet. Und St. Petrus sagt in den Geschichten der Apostel im 10. Kapitel V. 35.: Wer den HErrn fürchtet und recht thut, der ist ihm angenehm. Dagegen spricht Salomon in seinen Sprüchen im 28. Kapitel V. 26.: Wer sich auf sein Herz verläßt, ist ein Narr. Und der Prophet Jeremias im 17. Kapitel V. 5. sagt: Verflucht ist der Mann, der sich auf Menschen verläßt und hält Fleisch für seinen Arm, und mit seinem Herzen vom HErrn weicht. Also lehret uns Gottes Wort, daß wir alle in Gottesfurcht leben und in allem Thun auf Gott allein vertrauen, und für Vermessenheit uns hüten sollen; und daß uns also werde Glück und Segen folgen. Davon lehret auch das 2. Kapitel.

\section{Nehemia erlangt Gewalt und Briefe, Jerusalem zu bauen}
% Source scan 690; block decisions in selection.json

\quellseite{690}

\section{Summarium}

Ist eine feine Historia, wie Nehemias Erlaubniß bei dem Könige Arthasasta erlanget, daß er gen Jerusalem ziehen und die Mauern und Thore bauen möge, V. 5. Auf daß wir sehen, wenn man in der Furcht Gottes und mit einem ernstlichen Gebet etwas anfähet, wie Nehemias droben gethan hat, daß Glück dabei ist. Wie dagegen Unglück sein muß, wo man ohne Gottesfurcht und aus Vermessenheit etwas anfähet, wie doch gemeiniglich jedermann thut.

\section{Votum}

Regiere uns GOtt durch seinen heiligen Geist, daß wir nichts thun aus Vermessenheit, sondern daß wir in seiner Furcht alles anfangen und vollbringen, was ein jeder in seinem Beruf schuldig ist, und lasse die Werke unserer Hände wohl gerathen, Amen.

\chapter{Das 3. Capitel}

\section{Der Bau zu Jerusalem gehet an}

a (Gewaltigen) Die Armen müssen das Kreuz tragen, die Reichen geben nichts.
% Source scan 691; block decisions in selection.json

\quellseite{691}

\chapter{Das 4. Capitel}

\section{Vorrede}

Sollen die Feinde göttlichen Wortes und Namens gefällt und erlegt werden, so müssen wir, ein jeder in seinem Beruf und Stande, mit ernstem Seufzen in unserm Vortheil des lieben Gebets liegen und Gott zum Freunde haben und behalten. Ursache: unsre Feinde alle sind nichts anders, denn eine zornige Ruthe unsers Gottes und HErrn, welche er über uns verhänget um unserer Sünde willen, welche wir anders nicht brechen werden, denn mit Büßen und Beten. Derwegen wird recht gesagt: Wider die Strafen ist eine Arznei, die wahre christliche Buße, daß sie dennoch gelindert werden, wo nicht gar abgewendet; und das Gebet der Christen ist der Rüstzeug, Wehre und Waffen wider die Feinde, und für die Freunde der Christenheit. Denn wer mit rechtem Ernst betet, der streitet auf seiner Seite für Freunde, auf der andern Seite wider die Feinde der Christenheit. So helfe Gott, daß wir heute auch also leben und beten, daß es ihm zu Ehren, und zu Nutz und unsern Feinden zum Untergange gelange, Amen. Das von höret auch an das 4. Kapitel.
% Source scan 692; block decisions in selection.json

\quellseite{692}

\section{Der Bau gehet auch unter den Hindernissen fort}

\section{Summarium}

Ist eine Historia, in welcher man siehet, daß der Teufel den Juden sonderlich feind gewesen und ihnen durch ihre Nachbarn alles Unglück zugerichtet hat, daß sie mit einer Hand arbeiten, mit der andern sich der Feinde haben wehren müssen, V. 16. 18. In solcher Noth hält sich Nehemias zu dem Gebet und tröstet sein Volk, Gott werde bei ihnen sein und der Feinde Rath zu nichte machen. Diese Historia ist ein fein Fürbild unsers Lebens hie auf Erden, da hat man immerdar wider den Teufel, die Welt und unser eigen Fleisch zu fechten, soll anders die
% Source scan 693; block decisions in selection.json

\quellseite{693}
christliche Kirche durch reine Lehre und gottseeligen Wandel recht erbauet werden.

Helfe und regiere uns alle der treue GOtt, daß wir in fürstehenden Nöthen in täglicher Buße leben, in unserm Beruf treu sein und fleißig beten; und beschütze uns auch gnädiglich wider den Türken, und alle andere unsere Feinde, Amen.

\chapter{Das 5. Kapitel}

\section{Vorrede}

Der Wucher ist wider Gott und hat den zum höchsten Zorn und manchfaltigen Strafen bisher verursacht. Der Wucher ist eine Schande und böser Fürwurf unserm heiligen Evangelio. Der Wucher ist ein Hinderniß aller Ehrbarkeit und guter Polizei-Ordnung. Der Wucher ist ein Aergerniß aller frommen Herzen. Der Wucher ist eine Zerrüttung aller göttlichen und christlichen Nahrung. Der Wucher ist ein Gift der christlichen brüderlichen Liebe. Der Wucher ist ein Stifter aller Abgötterei. Der Wucher ist ein gewisses Pulver mit Schwefel anzuzünden, das ewig brennende Feuer im Gewissen. Wohl nun zeitlich, wohl nun ewig allen, die den Wucher unterlassen. Denn denen will Gott geben allerlei leiblichen und zeitlichen Segen, an Leib und Seele, Ehre und Gut, an ihren Kindern und Nachkommen. Ja es soll ihnen auch folgen der ewige himmlische Segen, ewiges Leben und Seligkeit, als bei denen, die ihren allein seligmachenden Glauben in rechtschaffenen Früchten der Liebe gegen den Nächsten sehen lassen. Wer Ohren hat zu hören, der höre. Das soll auch betrachtet werden bei dem 5. Kapitel.

\section{Nehemia thut die Beschwerden des Volks ab}
% Source scan 694; block decisions in selection.json

\quellseite{694}

\section{Summarium}

Ist eine feine Historia, darin Nehemias gebeut, daß man in der Theurung den armen Mann mit Wucher nicht übersezen, sondern was man ihnen abgewuchert habe, wieder geben und ihnen umsonst fürstrecken und leihen soll, um der Furcht Gottes willen, V. 11. Solches wäre heutiges Tages ein sehr nöthiges Exempel, daß die weltliche Obrigkeit auch also thäte und den Reichen wehrete, daß der arme Mann, sonderlich wenn theure Zeit ist, von ihnen nicht beschweret würde. Zum andern ist ein Exempel: nachdem Nehemias zum Landpfleger verordnet ist, er solches Amt also verwaltet, daß seine Leute mit unnöthigen Unkosten nicht beschweret werden, V. 14. Zum Exempel aller Obrigkeit, daß sie dergleichen mit ihren Unterthanen auch so handeln, und mehr auf derselben Nutz, Zunehmen und Gedeihen, denn auf den eigenen Nutz sehen. Aber man erfähret leider das Widerspiel.

\section{Votum}

Gott, unser lieber Herr und Vater, regiere unsere Obrigkeit, daß sie ihres Amts treulich pflege und der Unterthanen Wohlfahrt fleißig fördere; und helfe, daß alle Wucherer und Geizhälse heut Buße thun, der zeitlichen und ewigen Strafe zu entfliehen, Amen.

\chapter{Das 6. Kapitel}

\section{Vorrede}

Afterreden, den Nächsten verleumden, belügen und verrathen ist leider gemein worden und hat als eine pestilenzische Gift und vergifteter Wurm alle Winkel durchkrochen, und ist keine Kirche so wohl bestellet, keine Festung so wohl verwahret, keine Stadt, Dorf, noch Haushaltung dazu so wohl versehen, da dieser Wurm nicht eingeschlossen wäre. Derwegen, erstlich hüte man sich männiglich für solchem Laster und rechten Malzeichen des Teufels. Denn wer seinen Nächsten heimlich verleumdet, den vertilge ich, sagt Gott, der Herr, im 101. Psalm V. 5. Und St. Paulus spricht 1 Cor. 6, 24.: Die Lästerer werden das Reich Gottes nicht ererben. Dagegen sagt David im 15. Psalm: Herr, wer wird wohnen in deiner Hütte? Wer wird bleiben auf deinem heiligen Berge? Wer die Wahrheit von Herzen redet, und wer mit seiner Zunge nicht verleumdet. Und Esaias im 33. Kapitel V. 15. spricht: Wer in Gerechtigkeit wandelt, und redet, was recht ist. Wer Unrecht hasset, samt dem Geiz, der wird in der Höhe wohnen. Fürs andere, widerspreche auch ein jedes dem Lügen und Afterreden, wenn es solches von Verleumdern und Ehrendieben müßte hören. Denn, wie Salomon sagt in seinen Sprüchen im 17. Kapitel V. 15.: Wer den Gottlosen Recht spricht, und den Gerechten verdammt, die sind beide dem Herrn ein Greuel. Dies erinnert auch das 6. Kapitel.

\section[Nehemia treibet durch seine Vorsichtigkeit die Hindernisse am]{Nehemia treibet durch seine Vorsichtigkeit die Hindernisse am Werke zurück}
% Source scan 695; block decisions in selection.json

\quellseite{695}

\section{Summarium}

Dies ist eine Historia, wie die heidnischen Amtleute der Hoffart nach sich an Nehemia versuchen, ob sie ihn mit Lügen und anderm könnten verzagt machen und abschrecken, V. 2., bestellen auch etliche falsche Propheten wider ihn, V. 14. Solche Historia zeuget ehrlich, wie mancherlei Gefahr der fromme Nehemias auf sich hat müssen nehmen, darum, daß er treulich mit Jerusalem und den Juden gemeinet hat. Aber weil er Gott fürchtet und ihm vertraut, verachtet er solche Gefahr und bleibet fest bei dem, daß er angefangen hatte. Zum Exempel aller Obrigkeit, daß sie ihres Amts mit Fleiß aus warten soll, und gar nichts darnach fragen, es lache die Welt oder sehe sauer dazu, wenn es nur Gott gefällt und mit gutem Gewissen geschiehet.

\section{Votum}

Der treue Gott verleihe auch ferner unser lieben Obrigkeit starken Muth und andere nöthige Gaben zu glückseliger Regierung, und steure und wehre allen ihren Feinden, Amen.

\chapter{Das 7. Kapitel}

\section{Wächter der Stadt; Zahl und Geschenke des Volks}
% Source scan 696; block decisions in selection.json

% Scan 696: no apparatus retained; only biblical narrative/running matter.
% Source scan 697; block decisions in selection.json

\quellseite{697}

\chapter{Das 8. Kapitel}

\section{Vorrede}

Die zwei Testamente in der heiligen Biblia, sagt St. Hieronymus, sind die zwei Mutter-Brüste der Kirche, dadurch sie ihre KirchKinder säuget und ernähret. Derwegen ist das Alte Testament in hebräischer Sprache von den Propheten geschrieben worden, dieweil es der Juden Muttersprache und jedermann bei ihnen verständlich war; das Neue Testament aber ist in griechischer Sprache gestellet, dieweil dieselbe der Zeit die allgemeinste Landsprache und fast an denen Orten allen, in welchen und an welche die Bücher daselbst geschrieben worden, bräuchlich gewesen. Auch ist das Alte Testament in die griechische, und nachmals, sammt dem Neuen, von vielen Gelehrten, und sonderlich von St. Hieronymo in die lateinische und slavonische Sprache übersetzt worden, und darauf in die deutsche vom Herrn Luther, Herrn Philippo und ihren Mitgehülfen. So danken wir heute abermal herzlich Gott, unserm lieben Vater, daß wir die heilige Biblia haben, und daß uns diese alle Tage fürgelesen wird. Derselbe helfe gnädiglich, daß wir die heilige Biblia allzeit anhören, aus herzlicher Begierde, seinem Willen und die Wahrheit zu lernen, und mit ernstlichem Fürsatz, derselben zu glauben und darnach zu leben. Darauf höret ferner an das 8. Kapitel.

\section{Nehemia und Esra richten den Gottesdienst an}
% Source scan 698; block decisions in selection.json

\quellseite{698}

\section{Summarium}

Ist eine historia, wie Nehemias verordnet, daß man predigen und dem Volk das Gesetz fürlesen soll. Denn wo das Predigtamt nicht im Schwange gehet, da ists unmöglich, daß nicht Irrthum und falscher Gottesdienst sollte folgen. Wiederum, wo Gottes Wort ist, da werden die Herzen recht unterrichtet, daß sie ihre Sünde und Gottes Barmherzigkeit erkennen lernen, und thun nicht, was sie gut dünket, sondern was sie hören, das Gott zu thun befohlen hat. Wie sie hie das Laubhüttenfest und anders anrichten, V. 16. Um solche Wohlthat soll man Gott loben und danken, wie der Priester Esra hie auch thut.

\section{Votum}

Das ewige wesentliche Wort GOttes, unser HErr JEsus Christus, erhalte uns die heilige Biblia; und verleihe, daß die in unseren Kirchen bis an den jüngsten Tag, zu Stärkung des Glaubens, zu Besserung des Lebens, und zu Trost in Noth und Tod, allen Zuhörern vorgelesen werde, Amen.

\chapter{Das 9. Capitel}

\section{Vorrede}

Eines jeden Christen Leben soll eine tägliche Buße sein. Denn wir sündigen alle leider alle Tage. So ist auch die Stunde des Todes sehr unsicher und ungewiß, wenn, wie und wo eines von dieser Welt abscheiden soll; ja niemand hat nur auf eine Stunde Siegel und Brief, dieselbige zu überleben. So sagen sehr fein die alten Lehrer: Durch die Buße wird der Mensch entbunden von seinen Sünden; er wird der christlichen Kirche wieder eingeleibet, mit Gott versöhnet, in ihm werden die geistlichen Gaben gemehret, aus einem Teufels-Kind wird ein Gottes-Kind, das Paradies wird ihm geöffnet. Ja sie sagen ferner: Die Buße erfreut die Engel,
% Source scan 699; block decisions in selection.json

\quellseite{699}
sie macht gerecht den Gottlosen, sie bringt wieder die verlornen Güter, sie schenket das ewige Leben, sie behütet für der ewigen Strafe. Ach so wollen wir alle Tage unsere Sünde erkennen, um Vergebung in wahrem Glauben bitten und denn beständig durch Gottes Hülfe in Gottesfurcht wandeln und also in täglicher Buße leben. Davon lehret auch das 9. Kapitel.

\section{Des Volcks öffentliche Buße}

\section{Summarium}

Hie ist eine Historia, wie das Volck zusammen kommen, eine Fasten gehalten, V. 1., und sich auf ein neues mit Gott verbunden und die Leviten eine schöne lange Predigt gethan haben. In welcher sie für dem Volke ehrlich Gott loben und danken für alle seine Wohlthaten, die er sonderlich den Juden je und je hat widerfahren
% Source scan 700; block decisions in selection.json

\quellseite{700}
lassen. Darnach bekennen sie, wie sie sich gegen solchem gnädigen Gott (leider) allweg ungebührlich gehalten, seines Wortes nicht geachtet und seine Propheten erwürget haben, V. 16. Darum was ihnen für Unglück und Strafe begegnet sei, das haben sie alles wohl verdienet, V. 30. Auf solche Beichte folget hernach ein Gebet, daß Gott sie wieder zu Gnaden annehmen, alles Elend wenden und ihnen gnädiglich helfen wolle.

Auf diese Weise sollen wir auch beten lernen, daß wir Gottes Güte preisen, unsere Untugend und Sünde für ihm bekennen und beichten und weiter seiner Gnade warten, und uns bessern.

a Himmels Brod.
% Source scan 701; block decisions in selection.json

\quellseite{701}

\section{Summarium}

Mit eben dem vorhergehenden Summarium.

\section{Votum}

Helffe GOtt, daß wir alle, jung und alt, reich und arm, diese nöthige und selige Betekunst wohl lernen, täglich fleißig und fruchtbarlich üben, und GOtt auch dienen in allen Ständen, in Heiligkeit und Gerechtigkeit, Amen.

\chapter{Das 10. Kapitel}

\section{Vorrede}

Der HErr Christus verbindet sich gegen uns in der Schrift, daß er die Almosen reichlich bezahlen wolle, zeitlich, geistlich und ewig. Zeitlich will Gott die Almosen belohnen. Denn also sagt Christus, Lucas im 6. Kapitel V. 38.: Gebet, so wird euch gegeben. Ein voll, gedrückt, gerüttelt und überflüssig Maß wird man in euren Schooß geben; denn eben mit dem Maß, da ihr mit messet, wird man euch wieder messen. Geistlich will Gott die Almosen belohnen. Denn also spricht abermal Christus, Matth. 5, 7.: Selig sind die Barmherzigen, denn sie werden Barmherzigkeit erlangen. Ewiglich will Gott die Almosen belohnen. Denn also saget Christus wiederum Lucä 16, 9.: Machet euch Freunde mit dem ungerechten Mammon, auf daß, wenn ihr nun darbet, sie euch aufnehmen in die ewigen Hütten. Das sind wahrlich treffliche Zusagungen. Aber sehet, wie ein wunder Ding ist es. Wenn euch ein Mensch, der doch von Art ein Lügner ist, etwas zusaget, so könnt ihr ihm glauben, ihm leihen und borgen, der doch keine Stunde seines Lebens und Gutes sicher ist. Warum thut ihr solches auch nicht gegen eurem HErrn Christo? Der wahrhaftig und ewig ist, der nicht schlechte Zusagung gethan hat, sondern ist auch darauf gestorben zu genugsamer Versicherung. Sonderlich aber soll man zur Unterhaltung des Gottesdienstes gern und treulich helfen, und solches will euch auch der HErr Christus treulich vergelten, zeitlich, geistlich und ewig. Davon höret auch an das 10. Kapitel.

\section{Versiegelung des erneuerten Bundes}
% Source scan 702; block decisions in selection.json

\quellseite{702}

\section{Summarium}

Dies ist eine Historia, wie das Volk auf ein neues mit Gott sich verbunden habe, daß sie fort nicht mehr gottlos sein, sondern nach dem Willen Gottes leben wollen, V. 33. Und ist hie sonderlich zu merken, daß sie sich neben andern auch des verpflichteten, sie wollen jährlich ein jeder zum Tempel und Unterhaltung des Gottesdienstes ein Drittteil eines Orts vom Gulden geben, V. 29. Denn Geld und Gut kann man nicht besser und seliger anlegen, denn so man damit den rechten Gottesdienst fördert.

\section{Votum}

Der Vater aller Barmherzigkeit, unser lieber GOtt, mehre und erhalte allen ihre Güter, die solche zu seiner Ehre gebrauchen, und damit seine Gottesdienste erhalten helfen, und tröste und bewahre sie zum ewigen Leben in der letzten Noth, Amen.

\chapter{Das 11. Capitel}

\section[Register der neuen Einwohner in der Stadt und]{Register der neuen Einwohner in der Stadt und auf dem Felde}
% Source scan 703; block decisions in selection.json

\quellseite{703}

\chapter{Das 12. Capitel}

\section{Vorrede}

Im 148. Psalm V. 12. sagt der heilige Prophet und König David: Jünglinge und Jungfrauen, Alte mit den Jungen, sollen loben den Namen des HErrn. Zuvor auch im 117. Psalm V. 1. vermahnet er: Lobet den HErrn alle Heiden; preiset ihn alle Völker. Denn seine Gnade und Wahrheit waltet über uns in Ewigkeit. Halleluja. St. Paulus aber lehret an die Galater im 6. Kapitel V. 6. also: Der unterrichtet wird mit dem Wort, der theile mit allerlei Gutes dem, der ihn unterrichtet. Und Sirach im 7. Kapitel V. 33. saget auch also davon: Ehre den Priester, und gib ihnen ihr Theil, wie dir geboten ist. Derwegen soll sich ein jeder Christ mit höchstem Fleiß dankbar erzeigen, gegen Gott wegen seiner Wohlthaten und dann gegen den Dienern göttlichen Wortes, wegen ihrer Mühe, Arbeit und Fleiß, so sie an dieselbe wenden, zu ihrer Bekehrung und Seligkeit, in Christo Jesu unserm HErrn. Von diesen beiden erinnert auch das 12. Kapitel.

\section[Namen der Priester und Leviten. Einweihung der Stadt.]{Namen der Priester und Leviten. Einweihung der Stadt. Ordnung der heiligen Pfleger}
% Source scan 704; block decisions in selection.json

\quellseite{704}

a) Jaddua dies ist Jaddua, der dem großen Alexander zu Jerusalem empfing. Also hat Nehemia, vielleicht Esra auch, den großen Alexander erlebt, und ist sehr alt worden, und vier oder fünf Hoherpriester überlebt.
% Source scan 705; block decisions in selection.json

\quellseite{705}

\section{Summarium}

Ist eine historia, wie das Volk Gott danket und fröhlich ist, daß die Mauren an der Stadt wieder aufgebauet und Jerusalem ein wenig gebessert ist, V. 27.

Der Text meldet hie wiederum, wie man sich sonderlich darum habe angenommen, daß die Priester und Leviten mit der Nahrung versorget würden, daß sie der Kirchen desto fleißiger warten könnten, V. 24. Denn soll der Gottesdienst recht erhalten werden, so muß man etwas darauf wagen, daß tüchtige Personen erzogen und nach Nothdurft versorget werden.

\section{Votum}

Wir danken GOtt herzlich für alle Wohlthaten, die er unserer Stadt geistlich und leiblich erzeiget hat; der helfe, daß wir ihn und seine Diener ferner ehren, und daß solches auch alle unsere Nachkommen thun mögen, bis auf den jüngsten Tag, Amen.

\chapter{Das 13. Capitel}

\section{Vorrede}

Wie allezeit böse verkehrte Leute gewesen, die in der christlichen Kirche Aergerniß angerichtet haben: also findet man solche Leute noch in allen christlichen Gemeinen und Sammlungen, und wohl solche, die feindseliger in Haß und Widerwillen, als die Heiden, ärger in Unzucht, denn die Türken, tiefer im Wucher, denn die Juden leben. So sollen nun die Christen von der Kirche Gottes nicht abfallen oder sich absondern, solchem Aergerniß wegen, sondern sie sollen ihre Herzen verwahren mit rechter Lehre und fleißigem Gebet, und also alle Aergernisse überwinden und Gliedmaßen und Bürger der Kirche und Gemeine des Herrn Christi bleiben. Wie St. Paulus vermahnet 1 Cor. 16, 13.: Wachet, stehet im Glauben, seid männlich und stark. Und unser Herr Jesus Christus lehret dieses, da er Matth. 18, 8. also sagt: So deine Hand, oder dein Fuß dich ärgert, so haue ihn ab und wirf ihn von dir. Das ist, so dich irgend ein Mensch, er sei Freund oder Feind, durch seinen sündlichen Wandel, von Gott und von der Gottseligkeit abziehen will, so sollst du ihm nicht folgen, sondern widerstreben und fliehen. Denn durch das Wort Hand verstehet der Herr Christus unsere Freunde, welcher Freundschaft uns nicht weniger vonnöthen ist, als unsere rechte Hand. Fuß aber heißt so viel als der Wandel und Thun des Menschen. Darauf höret nun an das 13. und letzte Kapitel.

\section[Nehemias Eifer in Verbesserung der Policei, des Kirchen-]{Nehemias Eifer in Verbesserung der Policei, des Kirchen- und des häuslichen Wesens}
% Source scan 706; block decisions in selection.json

\quellseite{706}

\section{Summarium}

Das sind zwei Historien. Die erste, weil im Gesetze verboten war, daß kein Ammoniter sollte in die Gemeine Gottes kommen, will Nehemias nicht gestatten, daß Tobias, der Ammoniter, einen Kasten im Tempel habe, für sein und seines Hauses Opfer, wirft solches alles aus dem Tempel, lässet sichs nicht anfechten, daß der Gottesdienst desselbigen Gelds mangeln soll, sondern reiniget den Kasten.

Die andere Historie ist, daß die Juden bald nachließen, die Zehenten und Opfer zu geben, V. 10. Daraus konnte Nehemias leichtlich abnehmen, daß solches am Gottesdienst einen Nachtheil würde bringen. Denn wer seiner Nahrung warten und nachgehen muß, der kann der Kirchen nicht warten. Darum straft er solche Sünde, und verordnet von neuem, daß den Priestern und Leviten ihr Einkommen sollte gereicht werden, V. 11. Solches wäre sehr noth, daß es in allen Dörfern und Flecken von der Obrigkeit auch also bestellet würde.
% Source scan 707; block decisions in selection.json

\quellseite{707}

\section{Summarium}

Ist ein sehr fein Exempel, daß ein Regent allenthalben ein fleißig Aufsehen soll haben, daß die Leute nicht unrecht thun, und sich versündigen. Wie Nehemias mit dem Sabbath thut und mit denen, so heidnische Weiber nahmen.

Daß er nun bittet, Gott wolle an solche und dergleichen andere gute Werke denken, V. 22., ist recht gethan. Denn eben wie man um der Furcht Gottes willen das Uebel unterlassen soll: also soll man gute Werke darum thun, daß Gott verheißen hat, dieselbigen reichlich zu vergelten. Solche Hoffnung auf Gottes Verheißung gefällt Gott wohl, und hindert das Hauptstück christlicher Lehre gar nichts. Denn Nehemias bittet nicht, daß Gott um solcher guter Werke willen seiner wolle schonen, und ihn selig machen, sondern also spricht er: Herr, mein Gott, gedenke mir deß auch und schone mein, nach deiner Barmherzigkeit. Das heißt, nicht auf eigen Werk und Verdienst, sondern auf Gottes Güte und Gnade gebaut.

\section{Votum}

Verleihe Gott, der HErr, daß unsere liebe Obrigkeit auch ferner über der wahren Religion fleißig halte; und helfe, daß wir Zuhörer das Predigtamt auch treulich fördern, und segne, zeitlich und ewig, an Leib und Seele, uns alle, Amen.

\part{Das Buch Esther}

\chapter{Das 1. Kapitel}

\section{Vorrede}

Die Historia Esther ist eine schöne liebliche Geschichte, darinnen beide, hohe und niedere Personen, Obrigkeit und Unterthanen sich spiegeln können, und daraus lernen, wie sie sich für Unglück hüten, das Kreuz, so uns Gott aufgelegt, mit Geduld ohne Murren tragen, den Glauben ohne Zweifel mit ernstigem Gebet und Hoffnung zu Gott üben, und in christlichem untadeligen Leben fortfahren sollen. Auch ist darinnen schön und ganz artig abgemalet das Wesen und der Lauf der ganzen Welt, daß man sehen mag, wie es noch heutiges Tages in allen Ständen zugehe. Welches denn nicht allein nützlich, sondern auch nöthig ist zu wissen; sonderlich weil sich die Menschen mehr nach Exempeln, denn nach bloßen Worten, Strafen, Dräuen oder Vermahnungen richten. Solche Exempel aber sollen gebraucht werden zur Stärkung des Glaubens und Besserung des Lebens. Denn wer Gott vertrauet und auf ihn bauet, wird gleichermaßen, wie andere Gläubigen, von Gott aus allen Nöthen erlöset werden. Und dagegen werden

die, welche mit den Gottlosen in gleichen Sünden leben und verharren, mit gleicher Strafe auch heimgesucht werden. Denn das ist gewisser denn gewiß, daß der gerechte Gott keine Sünde ungestraft, und daß der barmherzige Gott kein gut Werk unbelohnet läßt, nach dem Gesetze zu reden. Dies ist die Summa des Buches Esther.

\section{Vasthi wird wegen verachteter Mahlzeit von Ahasveros verstoßen}
% Source scan 708; block decisions in selection.json

\quellseite{708}

\section{Summarium}

In diesem Buch ist eine Historia, wie Gott den Juden, so im Lande unter den Heiden waren, eine sonderliche Gnade bewiesen und sie für einem großen greulichen Jammer behütet habe.

Im 1. Capitel ist eine Historia, wie es aus Gottes Gericht wunderbärlich sich zugetragen habe, daß die ungehorsame Königin Vasthi entsetzet worden ist; auf daß Esther eingesetzt, und durch sie die Juden, so noch in Babel und sonst zerstreuet waren, errettet würden. Wie wir hernach in der Historia weiter werden hören.

Sonderlich aber ist hie zu merken, daß der Ungehorsam der Weiber gegen ihre Männer eine Sünde sei. Denn Gott hat dem Weibe um der Sünde willen die Strafe aufgelegt, daß sie unter des Mannes Gewalt leben und nicht ihr selbst Herr soll sein, wie Moses meldet, 1 Mos. 3, 9.
% Source scan 709; block decisions in selection.json

\quellseite{709}

\section{Votum}

Dem frommen lieben GOtt sei Lob und Dank gesagt, der seine Kirche wunderbarlich in Gefahr und Noth allzeit erhält; der verleihe, daß alle Eheweiber in schuldiger Liebe, Demuth und Gehorsam, gegen GOtt und ihren Ehemännern, beständig leben, Amen.

\chapter{Das 2. Kapitel}

\section{Vorrede}

Gott ist ein Vater der Armen, ein Tröster und Helfer der Waisen und der Wittwen-Richter. Daher sagt David im 146. Psalm V. 9.: Der Herr behütet Fremdlinge und Waisen, und erhält die Wittwen; und kehret zurück den Weg der Gottlosen. Darnach: Gott verhänget auch über Untreue und Verrätherei, daß sie offenbar werden und ihre verdiente Strafe endlich überkommen. Daher vermahnet der Prediger Salomonis im 10. Kapitel und sagt V. 20.: Fluche dem Könige nicht in deinem Herzen, und fluche dem Reichen nicht in deiner Schlafkammer, denn die Vögel des Himmels führen die Stimme und die Fittige haben, sagen's nach. Beides lehret auch das 2. Kapitel.

\section{Esther wird zu königlichen Ehren erhoben}
% Source scan 710; block decisions in selection.json

\quellseite{710}

\section{Summarium}

Ist eine feine Historia, wie es der frommen Esther glücklich gegangen und wohl gerathen sei, V. 8. Zum Exempel allen denen, so in Gottesfurcht wandeln und fromm sind. Denn ob sie wohl im Elend und allerlei Unglück sind, so will dennoch endlich Gott sie nicht lassen, sondern mit reichem Segen bei ihnen sein.

Am Ende ist eine Historia von den zweien Aufrührern, so den König Ahasveros ermorden wollten, und darum beide an Bäume gehängt worden, V. 21. Solch ihren Rath zeiget Mardachai Esther an, durch welche es wieder der König erfähret, V. 22. Solche Treue Mardachai wird hernach in der Historia ein wunderliches Urtheil verursachen.

\section{Votum}

Gott, der Herr, der das Gute treulich belohnet und das Böse ernstlich strafet, segne und behüte heut auch alle, die ihn fürchten, und steure allen Aufrührern, beide im geistlichen und weltlichen Regiment, und thue solches um der Ehre seines heiligen Namens willen, Amen.

\chapter{Das 3. Kapitel}

\section{Vorrede}

Gott hat die hohe Obrigkeit zum Ausschuß erwählt und ihre Gewaltigen ihnen zu Mittelsherren verordnet; doch daß sie alle ihre Anschläge und Fürnehmen, Thun und Lassen, nach seinem Wort, ihm zu Ehren und nach der christlichen Liebe ihren Unterthanen und Vertrauten zu Nutz, anstellen und verrichten sollen. Aber die fürnehmste Ursache, warum Gott Könige, Fürsten und andere Obrigkeit in der Welt verordnet hat, und warum er manchen so hoch anbringt, bei großen Potentaten und Herrn, ist, daß durch ihr Aufsehen, Fleiß und Sorge, Mühe und Arbeit, Friede auf Erden erhalten werde, daß nicht durch Krieg und Aufruhr, durch Rotten und Secten, christliche Religion zerstöret, Gottes Wort und Dienst niedergelegt, Zucht und Ehrbarkeit aufgehoben und die häusliche Nahrung verderbet oder verhindert werde. Aber viele mißbrauchen ihre Gewalt und Würde, auch wider Gott und seine christliche Gemeine, oder den Frommen und Unschuldigen zum Verderben und Untergang, so viel an ihnen ist. Als davon eine Historia meldet das 3. Kapitel.

\section{Hamans Ehrgeiz und Blutdürstigkeit}
% Source scan 711; block decisions in selection.json

\quellseite{711}

\section{Summarium}

Ist eine Historia, in welcher man ein sonderlich Bild siehet eines gottlosen Tyrannen. Haman kommt bei dem Könige zu großen Gnaden und Ehren und wird gewaltig, V. 1. Solche Gewalt macht ihn stolz und hoffärtig. Darum, daß ihn Mardachai nicht wollte anbeten, wie andere thaten, ward er ihm von Herzen gram, V. 5., und gedachte, wie die Tyrannen pflegen, seine Gewalt nicht allein wider ihn, sondern auch wider alle Juden zu üben, V. 7. Und bringets auch so weit, daß er bei dem Könige ein Gebot ausbringet, daß man alle Juden in allen Landen auf einen Tag erwürgen, ihr Gut rauben und Preis machen soll. V. 12. Da siehet man, wie der Teufel allewege Gottes Volk und der Kirchen so feind gewesen und die Tyrannen wider sie gereizet hat. Aber wie hernach folgen wird, will dennoch Gott dem Teufel seine Anschläge brechen und die Seinen erretten.

\section{Votum}

Gott, der Allmächtige, steure allen Tyrannen und tyrannischen Anschlägen und Fürnehmen wider seine Christenheit; und erhalte seine Kirche und deren Gliedmaßen, zu Nutz und Freude dem menschlichen Geschlechte, um Jesu Christi, unsers HErrn willen, Amen.

\chapter{Das 4. Kapitel}

\section{Vorrede}

Der heilige Apostel Johannes sagt 1. Epistel, 3. Kapitel V. 16.: Daran haben wir erkannt die Liebe, daß Christus sein Leben für uns gelassen hat, und wir sollen auch das Leben für die Brüder lassen. Und der HErr Christus sagt Matth. 10, 39.: Wer sein Leben findet, der wirds verlieren; und wer sein Leben verlieret um meinetwillen, der wirds finden. Hiermit werden wir gelehret und vermahnet, daß ein jeder viel lieber sein Leben in Gefahr stellen solle, denn auch einen einigen Christen in Angst und Noth stecken lassen. Deß ist ein Exempel Mosis, 2 Mos. 32, 32. Welcher begehrt, aus dem Buch des Lebens ausgetilget zu werden, wegen seiner Juden, denen er doch viel Guts gethan, und Gottes Zorn um öftern von ihnen abgewandt, da sie es doch nicht erkannt, sondern ihm viel Uebels erzeigt, also, daß sie ihn auf eine Zeit gedachten zu steinigen und zu tödten. Auch ist ein Exempel der Apostel Paulus, der da sagt Röm. 9, 3.: Ich begehre ein Anathema und Fluchopfer zu werden, daß meine Brüder, die Juden, möchten selig sein. Davon lehret auch das 4. Kapitel.

\section{Trauern und Fasten Mardachai und der Juden}
% Source scan 712; block decisions in selection.json

\quellseite{712}

\section{Summarium}

Ist eine solche, schöne Historig, wie Mardachai die Königin Esther vermahnet, daß sie ihr Leib und Leben in Gefahr setzen und bei dem Könige versuchen soll, ob das tyrannische Gebot wider die Juden möchte abgeschafft werden, V. 13.

Hier siehet man einen Unterschied zwischen Frommen und Gottlosen. Da Haman im Sinn hatte, etwas wider die Juden anzufangen, warf er das Loos und brauchte Zauberei; aber Esther und Mardachai, ehe sie anfangen, die Sachen zu handeln, demüthigen sie sich für Gott mit Fasten und ernstlichem Beten, darum gibt Gott auch Glück zu ihrer Handlung.

Am Ende ist ein fein Exempel, daß die fromme Esther ihr Leib und Leben in Gefahr setzet, ihr Volk zu erretten.

\section{Votum}

Helffe uns GOtt, daß wir auch einander herzlich meinen, und einander dienen und rathen, und derwegen keine Gefahr scheuen, und daß wir all unser Thun mit Gebet anfangen und glücklich verrichten, Amen.

\chapter{Das 5. Kapitel}

\section{Vorrede}

Der Königliche Prophet David spricht im 34. Psalm V. 7.: Da dieser Elende rief, hörte der HErr, und half ihm aus seinen Nöthen. Dieses wiederholet er und sagt V. 18.: Wenn die Gerechten schreien, so höret der HErr, und errettet sie aus all ihrer Noth. Das ist, das gläubige Gebet höret und erhöret Gott, und errettet, daß man nicht verderbe, aus aller Noth, wie groß und schwer, wie dick und breit die auch sein kann und mag. Darnach im 7. Psalm sagt aber David V. 15.: Siehe, der hat Böses im Sinn, mit Unglück ist er schwanger; er wird aber einen Fehl gebären. Er hat eine Grube gegraben und ausgeführt, und ist in die Grube gefallen, die er gemacht hat. Das ist, die Sünde ergreift zuletzt den Uebelthäter. Wer einem andern einen Stein in Weg legt, der fällt selbst darüber; wer einem andern einen Strick leget, der wird darinnen gefangen; wer untreulich redet, auf den wird es fallen, Spr. Sal. 27. Denn der gerechte Gott läßt nichts Böses ungestraft, und nichts Gutes unbelohnet, es stehe an lang oder kurz. Dieser Sprüche und Lehren Exempel ist das 5. Kapitel.
% Source scan 713; block decisions in selection.json

\quellseite{713}

\section[Esther gehet zum Könige, und Haman stellt Mardachai]{Esther gehet zum Könige, und Haman stellt Mardachai nach dem Leben}

\section{Summarium}

Ist eine schöne Historia, wie Gott der Esther Glück und Heil gibt, daß sie bei dem Könige in sondern Gnaden ist, V. 2. Denn was man mit Gott und dem Gebet anfänget, da will Gott mit Glück und Segen bei sein. Wiederum ists ein wunderlich Exempel mit Haman, daß alle Gottlosen schrecken und die Frommen wiederum trösten soll. Denn da ihm sein Unglück am nächsten ist, ist er am sichersten und zornigsten und hält für gewiß, ehe er morgen das Mittagmahl esse, soll Mardachai schon erhangen sein, V. 4. Aber Gott wendet’s wunderlich, und zum Trost, wenn wir fromm bleiben und an Gottes Wort hangen, so soll alles Unglück, das unsere Feinde uns zubereiten, über ihren eigenen Kopf kommen und wir errettet werden.

\section{Votum}

Der treue gütige GOtt erbarme sich auch unser in unserer Noth, erhöre unser Gebet und steure den Feinden seiner Kirchen und Christgläubigen Anschlägen und Werken, daß sie ohne Schaden abgehen und Er allezeit gelobet und gepreiset werde, Amen.

\chapter{Das 6. Kapitel}

\section{Vorrede}

Gott pfleget aus kleinen, verachteten Leuten große und ansehnliche zu machen, und aus großen kleine. Demgegen sollen wir lernen, daß die Hand und Gewalt Gottes stark genug sei, die Hoffärtigen zu stürzen und die Verdrückten zu erhöhen. Denn es ist keine Gewalt so groß, Gott kann die leichtlich zu nichte machen; auch ist keine Noth so groß, Gott kann aus solcher die Seinen erlösen. Denn die Hand des Herrn hilft gewaltiglich und kann alles ändern, sagt der königliche Prophet David, Ps. 20, 7. und 77, 17. Und der Herr ist gütig und eine Beste zur Zeit der Noth, und kennet die, so auf ihn trauen. Wenn die Fluth überher läuft, so macht ers mit derselben ein Ende, tröstet der heilige Prophet Nahum im 1. Kapitel V. 8. Davon höret auch an das 6. Kapitel.
% Source scan 714; block decisions in selection.json

\quellseite{714}

\section{Mardachai wird zu hohen Ehren erhoben}

\section{Summarium}

Sie ist eine Historia, in welcher man fürnehmlich sehet, daß Gott alle Sachen in seinen Händen hat und wunderbärlich alle Dinge ordnen und schicken kann. Mardachai hat keine Gedancken mehr, daß er der Wohltat genießen sollte, daß er den König gewarnet hatte, sondern stund in Gefahr Leibes und Lebens, sammt allen andern Juden. Dagegen war Haman sicher, und hielts gar für gewiß, Mardachai sollte in einer Stunde an einem Baum hangen. Aber Gott wendets zu beiden Theilen wunderbärlich, daß Haman selbst den Rath gibt und Mardachai dadurch zu den höchsten Ehren kommt, V. 10. Darum soll sich niemand der gottlosen Praktiken erschrecken lassen, sondern alle Hoffnung auf Gott stellen, als der unversehens alles ändern und bessern kann.

Daneben ist ein fein Exempel von dem Könige, daß er die Zeit nicht mit Spielen noch Zechen unnütz zubringet, da er nicht schlafen kann, wie doch an Höfen gebräuchlich, sondern läßt ihm zu Bett lesen die alten Historien, und höret zu, V. 1. Solches sollten alle Fürsten und Regenten merken, und dergleichen auch thun. Aber wie man sehet, man hat zu Hofe für dem Wild, für dem Zechen und Spielen, ich geschweige anderes dergleichen Thuns, nicht dißweil, daß man mit Büchern könnte umgehen. Und solches sündliches und ärgerliches Thun liebet den Herren dermaßen, daß sie ihres Amts darüber vergessen und andern befehlen, daß sie Amtes halber selbst thun sollten, Händel verhören, Uneinigkeit stillen, die unbillig Bedrängten schützen, jedermann Recht schaffen und was mehr der Obrigkeit zustehet.

\section{Votum}

Der König aller Könige, unser lieber Gott, helfe, daß alle Obrigkeit und ihre Amtverwalter recht thun; und siehe seiner Kirchen auch bei in diesen gefährlichen Zeiten und errete aus aller Noth alle, die Ihn anrufen, und auf Ihn hoffen und vertrauen, Amen.
% Source scan 715; block decisions in selection.json

\quellseite{715}

\chapter{Das 7. Kapitel}

\section{Vorrede}

Der königliche Prophet David sagt im 34. Psalm V. 20.: Der Gerechte muß viel leiden; aber der HErr hilft ihm aus dem allen. Und der heilige Hiob spricht im 20. Kapitel V. 5. 19.: Der Ruhm der Gottlosen stehet nicht lange, und die Freude des Heuchlers währet einen Augenblick? Wenn gleich seine Höhe in den Himmel reicht, und sein Haupt an die Wolken rühret; so wird er doch zuletzt umkommen wie ein Koth; allerhand Mühe wird über ihn kommen. Denn er hat unterdrückt und verlassen den Armen. Beides von Errettung der Frommen und Untergang der Gottlosen lehret David im 37. Psalm V. 5. also: Befehl dem HErrn deine Wege, und hoffe auf ihn, er wirds wohl machen, und wird deine Gerechtigkeit herfür bringen wie das Licht, und dein Recht wie den Mittag. V. 35.: Ich habe gesehen einen Gottlosen, der war trotzig, und breitete sich aus, und grünte wie ein Lorbeerbaum. Da man vorüber ging, siehe, da war er dahin; ich fragte nach ihm, da ward er nirgend funden. Dieser Sprüche Exempel an Mardachai und Haman erzählet das 7. Kapitel.

Haman wird an einen Baum gehänget.

\section{Summarium}

Ist eine tröstliche Historia für alle gottseelige Menschen, daß Gott nicht allein die Frommen retten und ihre Unschuld an Tag bringen, und sie ehren will, sondern auch die Tyrannen strafen und alles Unglück, das sie den Frommen gedacht, über ihren eigenen Kopf bringen. Mardachai, der fromme Jude, sollte an einen Baum erhangen werden; so schaffets Gott, daß man Haman, den Tyrannen, selbst daran hanget, V. 10. Solche Aenderung kann Gott mit den Tyrannen machen. Darum soll man Gott vertrauen und sich für den bösen Praktiken der Tyrannen nicht so sehr fürchten; Gott ist allewege mächtiger und weiser, denn sie.

\section{Votum}

GOtt, der HErr Zebaoth, steure und wehre heute auch allen Feinden seiner Christenheit, und tröste und erhalte alle, die seinen heiligen Namen lieben und ehren, bringe auch ihre Unschuld an den Tag, und stärke unsern Glauben, Amen.

\chapter{Das 8. Kapitel}

\section{Vorrede}

Der heilige Geist im 37. Psalm sagt: Ergrimme dich nicht über den, dem sein Muthwille glücklich fortgehet; denn die Bösen werden ausgerottet; die aber des HErrn harren, werden das Land erben. Es ist noch um ein Kleines, so ist der Gottlose nimmer; und wenn du nach seiner Stätte sehen wirst, wird er weg sein. Aber die Elenden werden das Land erben, und Lust haben in großem Frieden. Und wie ein Palmbaum die Eigenschaft hat, daß er mit seiner Last gebeuget, viel weniger unterdrückt werden kann; sondern je mehr man ihn beschweret, je höher er
% Source scan 716; block decisions in selection.json

\quellseite{716}
sich dagegen aufbäumet: also kann kein Fall die Schriften stürzen, wie groß er auch ist. Denn die Höllenpforten können sie nicht überwältigen, Matth. 16, 18. Also, wie das Gras daher grünet, daß einem Menschen das Herz im Leibe lacht, wenn ers so lieblich stehen siehet; aber bald wird abgehauen, da liegt es und verdorret, und wird von den Pferden und Kühen aufgefressen: also wendet sich auch bald das Blatt mit den Gottlosen, daß sie heute in Freude, morgen in Leid, heute im Leben, morgen im Tode und gestorben sind. Von diesen Wunderthaten Gottes lehret auch das 8. Kapitel.

\section[Die Juden haben Erlaubniß, sich an ihren Feinden]{Die Juden haben Erlaubniß, sich an ihren Feinden zu rächen}

a. Das ist: Weil die vorigen Briefe Hamans waren mit des Königs Siegel versiegelt, hätte es die Juden nichts geholfen, wo sie nicht wären von neuem durch andere Briefe widerrufen.

\section{Summarium}

Ist eine Historia, wie sich die Sachen gar verkehren, daß der König nicht allein das Gebot
% Source scan 717; block decisions in selection.json

\quellseite{717}
widerrufet, sondern auch den Juden erlaubet, auf den bestimmten Tag sich zu rächen an ihren Feinden, die ihnen übels gedacht, und sich ihres Unglücks erfreuet hätten, V. 13. Also will Gott allewege die Seinen schützen und die Tyrannen dämpfen, so fern man nur das Vertrauen auf ihn allein setzet; aber solche Hülfe verzeucht sich eine Zeitlang. Denn die Juden haben bis in die drei Monat in der Sorge und Gefahr stehen müssen, daß die Heiden auf einen Tag sie alle erschlagen und vertilgen, das Gut, Leib und Leben, Weib und Kind und alles rauben sollten.

\section{Votum}

Der treue barmherzige Gott wolle heute auch alle Tyrannen sammt ihrem Anhang stürzen; und schütze seine arme bedrängte Christenheit, und verleihe allen Kreuzträgern unter dem Verzug der Hülfe wahren Glauben und beständige Geduld, Amen.

\chapter{Das 9. Capitel}

\section{Vorrede}

Gott sagt im 89. Psalm V. 23.: Die Feinde sollen meinen Knecht David nicht überwältigen, und die Ungerechten sollten ihn nicht dämpfen; sondern ich will seine Widersacher schlagen vor ihm her, und die ihn hassen, will ich plagen: das ist, Gott will die Gläubigen und Frommen erhalten, stärken, wider alle Feinde beschützen, ihre Widerwärtige plagen und schlagen. Auch spricht David im 50. Psalm V. 14.: Opfere Gott Dank, und bezahle dem Höchsten deine Gelübde: das ist, erkenne o Mensch, und vergiß nicht, was Gott Gutes an uns gethan, daß er uns all unsere Sünde vergibet, unser Leben vom Verderben erlöst und uns mit Gnaden und Barmherzigkeit gekrönet hat. Dem sollt du auch mit Mund, Herzen und mit deinem ganzen Leben dafür danken, ihm vertrauen und ihn für den Brunnquell aller Güter erkennen und halten. Dieser, als des Trosts und Lehrspruchs Exempel ist das 9. Capitel.

\section[Die Juden stellen nach göttlicher Rache ein Freudenfest]{Die Juden stellen nach göttlicher Rache ein Freudenfest an}
% Source scan 718; block decisions in selection.json

\quellseite{718}

\section{Summarium}

Ist eine Historia, wie Esther und Mardachai die zween Tage, an welchen ihnen Gott gnädig geholfen und sie errettet, und an ihren Feinden gerochen hatte, verordneten, daß mans feiern soll, V. 18., und solches um der Ursache willen, daß man solcher großen Wohlthat nicht vergessen, sondern Gott dafür danken, und solches sich weiter in andern Nöthen erinnern und damit trösten möge, Gott werde auch helfen, wie er diesmal geholfen habe. Solcher Dank und Zuversicht auf Gottes Hülfe ist zuvor geboten in der ersten Tafel Mose. Daß man aber die Feiertage dazu verordnet, geschieht darum, daß der gemeine Mann solcher Ceremonien darf, sich damit auf einen gewissen Tag zu erinnern der vorigen Wohlthaten Gottes. Dergleichen Feste haben die Juden hernach mehr geordnet, als sie wieder ins Land sind kommen.

\section{Votum}

Gott, der hochgelobten Dreifaltigkeit, sei Lob, Ehr, Preis und Dank, der bisher auch seine Kirche wunderlich erhält und schützet; der vertilge ferner alle der Christenheit Feinde, und helfe uns, daß wir seine Gottesdienste allzeit in wahrem Glauben fleißig verrichten, Amen.

\chapter{Das 10. Capitel}

\section{Mardachai wird gerühmet}
